\begin{abstract}
% Spatially resolved transcriptomics (SRT) reveals gene expression patterns within the context of tissue microenvironment by in situ measuring gene expression at native spatial locations, which enables the deciphering of spatial domains. However, existing methods often simply model cellular dependencies using $k$-nearest neighbor graphs, overlooking the hierarchical characteristics and higher-order dependencies of cellular networks. To address this issue, we propose \method{}, a cross-view hypergraph-enhanced framework for identifying spatial domains in SRT data. \method{} captures contextual and hierarchical cellular dependencies using hypergraphs, performs hypergraph representation learning from both spatial- and expression-view, and facilitates cross-view knowledge transfer through a cross-attention mechanism and self-correlation reorganization. Additionally, a multi-view self-supervised clustering strategy is adopted to constrain \method{} to learn clustering-friendly representations, while hypergraph smoothing is applied to enhance consistency with biological priors. Finally, a zero-inflated negative binomial reconstruction is adopted to handle the inherent over-dispersion and dropouts in SRT data. Comprehensive experiments on multiple benchmark datasets demonstrate the effectiveness and superiority of our \method{} in various downstream tasks.
Single-cell RNA sequencing (scRNA-seq) serves a pivotal role in characterizing gene expression at the cellular level, enabling the identification of cell types and advancing the understanding of cellular heterogeneity. Despite the significant progress in scRNA-seq data clustering, we argue that current methods always ignore the sparsity and noise, as well as the complex intercellular structural information inherent in scRNA-seq data. Toward this end, in this paper, we propose a novel \underline{s}ingle-\underline{c}ell RNA-seq clustering framework via deep Siamese \underline{G}raph \underline{T}ransformer \underline{N}etwork (termed \method{}), which explicitly integrates gene expression profile and intercellular structural dependencies for cell clustering. In particular, we formulate scRNA-seq data as a graph and construct two augmented graph views that serve as dual views to capture complementary intercellular information. Then, a Siamese graph transformer network is employed to explicitly incorporate shortest-path information and node-wise distances for capturing richer structural relationships between cells. Finally, we employ an optimal transport strategy to guide the cell clustering in a self-supervised manner. Extensive experiments on multiple benchmark scRNA-seq datasets demonstrate that our \method{} consistently outperforms existing methods. Our code is available at \href{https://github.com/W-RMSL/scGTN}{https://github.com/W-RMSL/scGTN}.
\end{abstract}